% =============================================================================
% MLSys·im Tutorial Tutorial — Parts 0–4 (Morning Session)
% =============================================================================
\documentclass[aspectratio=169, 12pt]{beamer}
\usepackage{../../../slides/assets/beamerthememlsys}

\mlsyssetup{
  volume       = {Tutorial},
  chapter      = {Module 1},
  logo         = {../../../slides/assets/img/logo-mlsysbook.png},
  instlogo     = {../../../slides/assets/img/logo-harvard.png},
  chaptertitle = {MLSys·im: Foundations \& Architecture},
}

% --- Fonts ---
\usepackage[T1]{fontenc}
\usepackage[scaled=0.9]{helvet}
\usepackage{courier}
\renewcommand{\familydefault}{\sfdefault}

% --- Packages ---
\usepackage{amsmath}
\usepackage{booktabs}
\usepackage[table]{xcolor}
\usepackage{listings}
\usepackage{tikz}
\usetikzlibrary{arrows.meta, positioning, calc, decorations.pathreplacing}

% --- Code listings ---
\lstset{
  language=Python,
  basicstyle=\ttfamily\scriptsize,
  keywordstyle=\color{crimson}\bfseries,
  stringstyle=\color{datastroke},
  commentstyle=\color{midgray}\itshape,
  backgroundcolor=\color{computeblue!20},
  frame=single,
  rulecolor=\color{computestroke},
  numbers=none,
  breaklines=true,
  columns=fullflexible,
  keepspaces=true,
  showstringspaces=false,
  xleftmargin=4pt,
  xrightmargin=4pt,
  aboveskip=6pt,
  belowskip=4pt,
}

% --- Convenience macros ---
\newcommand{\mlsysim}{\texttt{mlsysim}}
\newcommand{\wallbox}[2]{%
  \begin{block}{#1}#2\end{block}%
}
\newcommand{\PredictStart}{\begin{alertblock}{Predict Before You Peek}}
\newcommand{\PredictEnd}{\end{alertblock}}

% --- Image paths ---
\graphicspath{{images/}}

% --- Section count (must match actual \section{} count) ---
\setcounter{mlsystotalsections}{6}

\title{MLSys·im Tutorial --- Module 1}
\subtitle{Foundations \& Architecture}
\author{Vijay Janapa Reddi}
\institute{Harvard University}
\date{Tutorial}

% =============================================================================
\begin{document}

% --- Roadmap ---
\begin{frame}[shrink]{Roadmap: Conference Tutorial}
\centering\small
\begin{tabular}{rll}
\toprule
\textbf{Module} & \textbf{Topic} & \textbf{Status} \\
\midrule
\rowcolor{crimson!12}
\textbf{Module 1} & \textbf{Foundations \& Architecture} & \textbf{$\leftarrow$ You are here} \\
Module 2 & Advanced Single-Node Analysis & \\
Module 3 & Scale, Dollars, and Carbon & \\
Module 4 & Design Space Exploration \& Synthesis & \\
\bottomrule
\end{tabular}
\end{frame}


\section{Welcome \& Setup}

% --- Slide 0.1: Title ---
\begin{frame}[shrink]
\note{[1 min] Welcome attendees, set the tone.
Welcome to the MLSys-im tutorial. Today we will build quantitative
intuition for ML systems from first principles.
% -- FLEX: [CORE] Title slide --- do not skip.
}
\titlepage
\end{frame}

% --- Slide 0.1b: The $200M Question ---
\begin{frame}[shrink]{The \$200 Million Question}
\note{[3 min] THE HOOK. Open strong. Don't touch your laptop. Look at the audience.}

\centering
\Large
\textbf{Meta spent \$200M training Llama-3-405B.}\\[1cm]

\normalsize
Before a single GPU was purchased:\\[0.3cm]
\begin{itemize}
  \item How would you know \textbf{16,384 H100s} was the right fleet?
  \item How would you know \textbf{405B parameters} was the right model size?
  \item How would you know it would take \textbf{54 days}, not 540?
\end{itemize}

\vfill
\small\textcolor{gray}{We will answer all three questions today --- on your laptop, in under a second, with no GPU.}
\end{frame}

% --- Slide 0.1c: Live Demo Reveal ---
\begin{frame}[fragile,shrink=10]{Answer in 0.1 Seconds}
\note{[2 min] Run this LIVE. The room should gasp at how fast the answer appears.}

\begin{lstlisting}
import mlsysim

profile = mlsysim.Engine.solve(
    mlsysim.Models.Language.Llama3_8B,
    mlsysim.Hardware.Cloud.H100,
    batch_size=1,
)
print(f"Bottleneck: {profile.bottleneck}")  # Memory
print(f"MFU: {profile.mfu:.3f}")            # 0.003
\end{lstlisting}

\vfill
\centering
\textbf{That took 0.1 seconds. On a laptop. No GPU.}\\[0.2cm]
\small Now imagine doing this for every hardware option, every model size,\\
every parallelism strategy, every region. \textbf{That is mlsysim.}
\end{frame}

% --- Slide 0.2: What You Will Learn Today ---
\begin{frame}[shrink]{What You Will Learn Today}
\note{[2 min] Walk through objectives quickly. Emphasize that by the
end of the day every attendee will be able to do these five things.
% -- FLEX: [CORE]
}

\small
By the end of this tutorial you will be able to:
\begin{enumerate}
  \item \textbf{Identify} which physical constraint is the binding bottleneck
        for any ML workload on any hardware.
  \item \textbf{Decompose} training and inference time using the Iron Law.
  \item \textbf{Compare} hardware configurations quantitatively with \mlsysim.
  \item \textbf{Reason} about the compute--memory--communication tradeoff space.
  \item \textbf{Estimate} TCO and carbon footprint for a real deployment.
\end{enumerate}

\vspace{0.3cm}
\centering
\textit{All you need is a laptop and} \texttt{pip install mlsysim}
\end{frame}

% --- Slide 0.3: Setup Check ---
\begin{frame}[fragile,shrink=10]{Setup: Install \& Verify}
\note{[3 min] Give attendees 2 minutes to run these commands.
Walk around and help anyone with pip issues.
If someone cannot install, they can pair with a neighbor.
% -- FLEX: [CORE] --- must verify before proceeding.
}

\small
Open a terminal and run:

\begin{lstlisting}
pip install mlsysim
python3 -c "import mlsysim; print(mlsysim.__version__)"
# Expected output: 0.1.0
\end{lstlisting}

\vspace{0.3cm}
Then run the hello-world sanity check:

\begin{lstlisting}
import mlsysim
model = mlsysim.Models.Language.Llama3_8B
hw    = mlsysim.Hardware.Cloud.H100
prof  = mlsysim.Engine.solve(model, hw, batch_size=1)
print(prof.bottleneck)   # -> "Memory"
\end{lstlisting}

\vspace{0.2cm}
\centering
\alert{If you see \texttt{Memory}, you are ready.}

\vspace{0.3cm}
\scriptsize
\textit{Convention for the rest of the day:}\\
\texttt{import mlsysim} is assumed.
We use \texttt{llama} $=$ \texttt{mlsysim.Models.Language.Llama3\_8B}
and \texttt{hw} $=$ \texttt{mlsysim.Hardware.Cloud.H100} as shorthands.
\end{frame}

% --- Slide 0.4: The 22-Wall Taxonomy ---
\begin{frame}[shrink]{The 22 Physical Walls of ML Systems}
\note{[2 min] This is the road map for the day. Point out that
we will hit walls 1--7 (Node) before lunch and walls 8--22
after lunch. Each wall has one equation and one mlsysim solver.
Ask: ``How many of these walls have you personally hit?'' Show of hands.
% -- FLEX: [CORE]
}

\scriptsize
\begin{columns}[T]
\begin{column}{0.48\textwidth}
\textbf{Domain 1: Node}\\[2pt]
\begin{enumerate}\setlength\itemsep{1pt}
  \item Compute Wall
  \item Memory Wall
  \item Software Wall (MFU)
  \item Serving Wall
  \item Batching Wall (KV cache)
  \item Streaming Wall
  \item Tail Latency Wall
\end{enumerate}

\vspace{0.2cm}
\textbf{Domain 2: Data}\\[2pt]
\begin{enumerate}\setlength\itemsep{1pt}\setcounter{enumi}{7}
  \item Ingestion Wall
  \item Transformation Wall
  \item Locality Wall
\end{enumerate}
\end{column}

\begin{column}{0.48\textwidth}
\textbf{Domain 3: Algorithm}\\[2pt]
\begin{enumerate}\setlength\itemsep{1pt}\setcounter{enumi}{10}
  \item Complexity Wall (Chinchilla)
  \item Reasoning Wall
  \item Fidelity Wall (Compression)
\end{enumerate}

\vspace{0.2cm}
\textbf{Domain 4: Fleet}\\[2pt]
\begin{enumerate}\setlength\itemsep{1pt}\setcounter{enumi}{13}
  \item Communication Wall
  \item Fragility Wall
  \item Multi-Tenant Wall
\end{enumerate}

\vspace{0.2cm}
\textbf{Domain 5: Operations}\\[2pt]
\begin{enumerate}\setlength\itemsep{1pt}\setcounter{enumi}{16}
  \item Capital Wall (TCO)
  \item Sustainability Wall
  \item Checkpoint Wall
  \item Safety Wall
\end{enumerate}
\end{column}
\end{columns}
\end{frame}

% --- Slide 0.5: The Iron Law (Preview) ---
\mlsysfocus{The Iron Law of ML Systems}{%
\[
T = \frac{\text{FLOPs}}{N \;\times\; \text{Peak} \;\times\; \text{MFU} \;\times\; \eta_{\text{scaling}} \;\times\; \text{Goodput}}
\]
\\[0.5cm]
\normalsize
Every wall maps to one of these five denominator terms.\\
This single equation is our compass for the entire day.
}

% =============================================================================
% RELATED WORK & POSITIONING (8 slides)
% =============================================================================
\input{related_work}

% --- Roadmap: You Are Here (Morning) ---
\begin{frame}[shrink]{Roadmap: You Are Here}
\note{[1 min] Quick orientation. We just finished the setup. Now the real work begins.}

\centering\small
\begin{tabular}{rll}
\toprule
\textbf{Time} & \textbf{Part} & \textbf{Status} \\
\midrule
9:00--9:30   & Part 0: Welcome \& Setup     & \checkmark Done \\
\rowcolor{crimson!12}
9:30--10:30  & \textbf{Part 1: Iron Law \& Roofline} & \textbf{$\leftarrow$ You are here} \\
10:45--11:45 & Part 2: Memory Walls \& Serving & \\
11:45--12:00 & Part 3: Compression & \\
\midrule
\textit{12:00--1:00} & \textit{Lunch} & \\
\midrule
1:00--2:15   & Part 4: Going Distributed & \\
2:30--3:15   & Part 5: Economics \& Sustainability & \\
3:15--3:45   & Part 6: Design Space Exploration & \\
3:45--4:15   & Part 7: TinyML to Frontier & \\
4:15--4:45   & Part 8: Advanced Topics & \\
4:45--5:00   & Part 9: Wrap-Up \& Capstone & \\
\bottomrule
\end{tabular}
\end{frame}

% =============================================================================
% PART 1: THE IRON LAW & ROOFLINE (15 slides)
% =============================================================================


\begin{frame}[shrink]{Constraint-Based Systems Engineering}
\textbf{The Concept:} Think of an ML system as a series of connected pipes.
\begin{itemize}
  \item \textbf{Data (Water):} Flows from storage $\to$ CPU $\to$ PCIe $\to$ GPU HBM $\to$ SRAM $\to$ ALU.
  \item \textbf{The Rule:} The system is exactly as fast as the \emph{narrowest pipe}.
\end{itemize}
\vfill
\begin{alertblock}{The Engineer's Goal}
The goal is \textbf{not} to optimize every component. The goal is to mathematically identify the \emph{binding constraint} (the narrowest pipe) and widen it. Optimizing a non-binding constraint yields exactly 0\% speedup.
\end{alertblock}
\end{frame}

\section{Architecture \& Registries}

\begin{frame}[shrink]{The \mlsysim{} Philosophy: Demand vs. Supply}
\textbf{Strict Separation of Concerns:}
\begin{itemize}
  \item \textbf{Demand (Workloads):} How many FLOPs? How many bytes of weights? (Abstract mathematical operations).
  \item \textbf{Supply (Hardware):} What is the peak TFLOP/s? What is the HBM bandwidth? (Physical silicon constraints).
\end{itemize}
\vfill
\begin{center}
\textbf{Solvers} sit in the middle, evaluating the intersection to find the \emph{binding constraint}.
\end{center}
\end{frame}

\begin{frame}[fragile,shrink=10]{The Type System \& Registries}
\mlsysim{} uses a strict, typed registry system powered by Pydantic. No magic dictionaries.
\begin{lstlisting}[language=Python]
from mlsysim.hardware.types import HardwareNode, ComputeCore, MemoryHierarchy
from mlsysim.core.constants import TFLOPs, GB, TB, second, watt

# Defining an accelerator using strict physical quantities
speculative_gpu = HardwareNode(
    name="Speculative GPU",
    release_year=2027,
    compute=ComputeCore(peak_flops=1200 * TFLOPs / second),
    memory=MemoryHierarchy(capacity=144 * GB, bandwidth=5.0 * TB / second),
    tdp=800 * watt
)
\end{lstlisting}
\end{frame}

\begin{frame}[shrink]{Zero Hallucinations: The Provenance Audit}
Academic tools must be reproducible. Every vetted number in \mlsysim{} is bound to a \textbf{Provenance Record}.
\begin{itemize}
  \item \textbf{Datasheets:} Must include a verified URL to the vendor specification.
  \item \textbf{Literature:} Must cite the published academic paper (e.g., Llama 3 networking topologies).
  \item \textbf{Estimates:} Must explicitly document the napkin-math assumptions.
\end{itemize}
\vfill
\begin{center}
\texttt{make audit} statically parses the AST to ensure \textbf{zero undocumented magic numbers}.
\end{center}
\end{frame}

\begin{frame}[fragile,shrink=10]{Dimensional Strictness (\texttt{pint})}
\begin{columns}[T]
\column{0.5\textwidth}
\textbf{The Problem:}
\begin{lstlisting}[language=Python]
# Is this MB/s or GB/s? 
# Wait, bits or bytes?
bandwidth = 400 
\end{lstlisting}

\column{0.5\textwidth}
\textbf{The \mlsysim{} Solution:}
\begin{lstlisting}[language=Python]
# Will throw runtime error if 
# divided by seconds instead of bits
bw = Q_("400 Gbps") 
speed = bw.to("GB/s") # 50 GB/s
\end{lstlisting}
\end{columns}
\vspace{1em}
Every API boundary strictly enforces SI units at runtime. This prevents silent mismatches when mixing networking (bits) and memory (bytes).
\end{frame}

\section{Iron Law \& Roofline}

% --- Slide 1.1: Key Question ---
\begin{frame}[shrink]{Key Question}
\note{[1 min] Pose the question dramatically. Pause for 5 seconds.
``This is the most important question in ML systems engineering.
By the end of this section you will answer it in 3 lines of Python.''
% -- FLEX: [CORE]
}

\centering
\Large\bfseries
Why doesn't doubling FLOPS\\[0.3cm]
double your throughput?
\end{frame}

% --- Slide 1.2: Constraints Drive Architecture ---
\begin{frame}[shrink]{Constraints Drive Architecture}
\note{[2 min] ``You don't choose a Transformer because it's trendy;
you choose it because of how it parallelizes on real silicon.
AI is not magic --- it is infrastructure, and infrastructure has laws.''
% -- FLEX: [CORE]
}

\small
\begin{itemize}
  \item Hardware has \textbf{finite compute} (FLOPS), \textbf{finite bandwidth}
        (GB/s), and \textbf{finite memory} (GB).
  \item Every workload demands some amount of each.
  \item The \textbf{binding constraint} is the one that takes the longest.
  \item \alert{You optimize the bottleneck, not the fast part.}
\end{itemize}

\vspace{0.5cm}
\centering
\begin{tikzpicture}[>=Stealth, node distance=3cm]
  \node[draw, fill=computeblue, rounded corners, minimum width=2.5cm, minimum height=1cm]
    (compute) {\textbf{Compute}};
  \node[draw, fill=datagreen, rounded corners, minimum width=2.5cm, minimum height=1cm,
    right=of compute] (memory) {\textbf{Memory BW}};
  \node[draw, fill=routingorange, rounded corners, minimum width=2.5cm, minimum height=1cm,
    right=of memory] (network) {\textbf{Network}};
  \draw[->, thick, crimson] (compute) -- node[above, font=\scriptsize] {which is slowest?} (memory);
  \draw[->, thick, crimson] (memory) -- node[above, font=\scriptsize] {which is slowest?} (network);
\end{tikzpicture}
\end{frame}

% --- Slide 1.3: The Roofline Model ---
\begin{frame}[shrink]{The Roofline Model (Williams et al., 2009)}
\note{[3 min] Draw the two regimes on the board. Left = memory-bound,
right = compute-bound. The ridge point is where they cross.
``Before I show numbers: if a model does 16B FLOPs and
loads 16 GB of weights, what is its arithmetic intensity?''
Expected: 1 FLOP/byte. That is far left on the Roofline.
WARN: Students conflate FLOPS with throughput.
% -- FLEX: [CORE]
}

\small
\[
\text{Attainable FLOPS} = \min\!\bigl(\text{Peak FLOPS},\;\;
\text{BW} \times \text{Arithmetic Intensity}\bigr)
\]

\vspace{0.2cm}
\begin{columns}[T]
\begin{column}{0.55\textwidth}
\textbf{Arithmetic Intensity} (AI):
\[
  \text{AI} = \frac{\text{FLOPs}}{\text{Bytes moved}}
  \;\;\bigl[\text{FLOP/byte}\bigr]
\]

\vspace{0.2cm}
\begin{itemize}\setlength\itemsep{2pt}
  \item \textbf{AI $<$ Ridge Point} $\Rightarrow$ \colorbox{datagreen}{Memory-bound}
  \item \textbf{AI $>$ Ridge Point} $\Rightarrow$ \colorbox{computeblue}{Compute-bound}
\end{itemize}

\vspace{0.2cm}
Ridge Point $=$ Peak FLOPS $/$ Peak BW
\end{column}

\begin{column}{0.42\textwidth}
\centering
\begin{tikzpicture}[scale=0.7]
  % axes
  \draw[->, thick] (0,0) -- (6.5,0) node[right, font=\scriptsize] {AI (FLOP/B)};
  \draw[->, thick] (0,0) -- (0,4.5) node[above, font=\scriptsize] {GFLOPS};
  % memory roof
  \draw[very thick, datastroke] (0,0) -- (3,3);
  % compute roof
  \draw[very thick, computestroke] (3,3) -- (6.2,3);
  % ridge point
  \fill[crimson] (3,3) circle (3pt);
  \node[above right, font=\scriptsize, crimson] at (3,3) {Ridge};
  % labels
  \node[font=\scriptsize, datastroke, rotate=42] at (1.2,1.7) {BW-limited};
  \node[font=\scriptsize, computestroke] at (4.8,3.4) {Compute-limited};
\end{tikzpicture}
\end{column}
\end{columns}
\end{frame}

% --- Slide 1.4: The Compute Wall ---
\begin{frame}[shrink]{Wall 1: The Compute Wall}
\note{[2 min] ``This is the speed limit. No software trick can make
your model run faster than the chip can crunch numbers.''
% -- FLEX: [CORE]
}

\small
\wallbox{The Compute Wall}{
\[
  T_{\text{compute}} = \frac{\text{Operations}}{\text{Peak FLOPS} \times \text{Efficiency}}
\]
}

\vspace{0.2cm}
\textbf{Example:} ResNet-50 inference at batch 256 on H100

\begin{itemize}
  \item FLOPs = $8.0 \times 10^{9} \times 256 = 2.05 \times 10^{12}$
  \item H100 FP16 Peak = 989 TFLOPS
  \item At 50\% MFU: $T = \frac{2.05 \times 10^{12}}{989 \times 10^{12} \times 0.5} \approx 4.1\;\text{ms}$
\end{itemize}

\vspace{0.2cm}
\alert{The chip is the ceiling. MFU is how close you get to it.}
\end{frame}

% --- Slide 1.5: The Memory Wall ---
\begin{frame}[shrink]{Wall 2: The Memory Wall}
\note{[2 min] ``Quick mental math: 16 GB model, 3.35 TB/s bandwidth.
How long to load?'' Give 10 seconds. Expected: 16/3350 = 4.8 ms.
WARN: Students assume compute is always the bottleneck because
GPUs are marketed on TFLOPS.
% -- FLEX: [CORE]
}

\small
\wallbox{The Memory Wall}{
\[
  T_{\text{memory}} = \frac{\text{Weight Bytes}}{\text{Memory Bandwidth}}
\]
}

\vspace{0.2cm}
\textbf{Example:} Llama-3 8B at batch size 1 on H100

\begin{itemize}
  \item Weight size (FP16) = $8\text{B} \times 2\;\text{bytes} = 16\;\text{GB}$
  \item H100 HBM3 BW = 3.35 TB/s
  \item $T = \frac{16}{3350} \approx 4.8\;\text{ms}$ just to load weights
  \item Meanwhile, compute finishes in $\sim$0.03 ms
\end{itemize}

\vspace{0.2cm}
\alert{At batch size 1, LLM inference is $\sim$\,160$\times$ memory-bound.}
\end{frame}

% --- Slide 1.6: Predict Before You Peek #1 ---
\begin{frame}[shrink]{Predict: H100 vs MI300X vs Gaudi\,3}
\note{[3 min] PREDICTION. Give the audience 60 seconds to think.
Expected answer: all memory-bound. BW ratios determine speedup, not FLOPS.
After the reveal, hammer home: ``The bottleneck determines the speedup.''
Turn to your neighbor: did you get the same answer? Why or why not? 60 seconds.
% -- FLEX: [CORE]
}

\centering
\Large\bfseries
Three flagship accelerators. Same workload.\\[0.3cm]
\normalsize
Llama-3 8B, batch size 1, FP16 inference.\\
Which is fastest---and by how much?\\[0.3cm]

\pause

\scriptsize
\begin{tabular}{lccc}
\toprule
              & \textbf{H100 (NVIDIA)} & \textbf{MI300X (AMD)} & \textbf{Gaudi\,3 (Intel)} \\
\midrule
Peak FP16     & 989 TFLOPS    & 1,307 TFLOPS  & 1,835 TFLOPS \\
HBM BW        & 3.35 TB/s     & 5.3 TB/s      & 3.7 TB/s \\
HBM Capacity  & 80 GB         & 192 GB        & 128 GB \\
\midrule
Bottleneck    & Memory        & Memory        & Memory \\
Weight-load time & 4.8 ms     & 3.0 ms        & 4.3 ms \\
\textbf{Speedup vs H100} & ---  & \textbf{1.6$\times$} & \textbf{1.1$\times$} \\
\bottomrule
\end{tabular}

\vspace{0.3cm}
\pause
\small
\alert{MI300X has fewer FLOPS than Gaudi\,3 but wins on bandwidth.\\
FLOPS don't determine speed when memory-bound.}
\end{frame}

% --- Slide 1.7: Live Demo --- Engine.solve (multi-vendor) ---
\begin{frame}[fragile,shrink=10]{Live Demo: Three Vendors, One API}
\note{[3 min] Run this live. The key moment: all three are memory-bound.
The ranking follows bandwidth, not FLOPS. This is ISCA---show that
mlsysim is not an NVIDIA-only tool.
% -- FLEX: [CORE]
}

\small
Run this in your Python session:

\begin{lstlisting}
import mlsysim

model = mlsysim.Models.Language.Llama3_8B
for hw_name in ["H100", "MI300X", "Gaudi3"]:
    hw = getattr(mlsysim.Hardware.Cloud, hw_name)
    p = mlsysim.Engine.solve(model, hw, batch_size=1)
    print(f"{hw.name}: {p.bottleneck}, "
          f"{p.latency:.2f}")
\end{lstlisting}

\vspace{0.2cm}
\begin{exampleblock}{What to look for}
  \begin{itemize}
    \item \texttt{bottleneck}: Memory for \textbf{all three}
    \item Ranking follows BW (MI300X $>$ Gaudi\,3 $>$ H100), not FLOPS
    \item Same API, same physics, different silicon
  \end{itemize}
\end{exampleblock}
\end{frame}

% --- Slide 1.8: MFU --- The Software Wall ---
\begin{frame}[shrink]{Wall 3: MFU --- The Software Wall}
\note{[2 min] ``MFU measures the gap between what the hardware could do
and what your software actually achieves. A 50\% MFU means you are
paying for twice the hardware you are using.''
% -- FLEX: [CORE]
}

\small
\wallbox{Model FLOPs Utilization}{
\[
  \text{MFU} = \frac{\text{Achieved FLOPS}}{\text{Peak FLOPS}}
\]
}

\vspace{0.2cm}
\begin{columns}[T]
\begin{column}{0.55\textwidth}
\textbf{What eats MFU?}
\begin{itemize}\setlength\itemsep{2pt}
  \item Kernel launch overhead
  \item Memory stalls (cache misses)
  \item Framework overhead (Python $\to$ CUDA)
  \item Suboptimal operator fusion
  \item \alert{Being memory-bound} (the biggest one!)
\end{itemize}
\end{column}
\begin{column}{0.42\textwidth}
\centering
\textbf{Typical MFU ranges}\\[4pt]
\scriptsize
\begin{tabular}{lc}
\toprule
Workload & MFU \\
\midrule
LLM training (optimized) & 40--55\% \\
LLM inference (bs=1) & $<$5\% \\
ResNet training & 30--40\% \\
FlashAttention & 60--75\% \\
\bottomrule
\end{tabular}
\end{column}
\end{columns}

\vspace{0.2cm}
\alert{Improving MFU is often cheaper than buying more GPUs.}
\end{frame}

% --- Slide 1.8b: What Is Eta? ---
\begin{frame}[shrink]{What Is $\eta$? (The Efficiency Parameter)}
\note{[3 min] CRITICAL — every demo uses eta. Explain it ONCE here, then it's understood for the day.
ANALOGY: ``eta is to ML systems what CPI is to CPU design --- an empirical constant that bridges peak specs and reality.''}

\small
\textbf{$\eta$ = Achieved FLOPS / Peak FLOPS} (Model FLOPs Utilization)

\vspace{0.3cm}
The gap between what your hardware \emph{could} do and what it \emph{actually} does.

\vspace{0.3cm}
\begin{columns}[T]
\begin{column}{0.48\textwidth}
\textbf{What reduces $\eta$:}
\begin{itemize}\setlength\itemsep{1pt}
\item Kernel launch overhead
\item SM occupancy limits
\item Memory coalescing misses
\item Framework overhead (Python GIL)
\item Communication stalls
\end{itemize}
\end{column}
\begin{column}{0.48\textwidth}
\textbf{Typical values:}

\scriptsize
\begin{tabular}{@{}lr@{}}
\toprule
Scenario & $\eta$ \\
\midrule
Training (Megatron-LM) & 0.40--0.55 \\
Training (PyTorch eager) & 0.08--0.15 \\
Inference decode, bs=1 & 0.01--0.05 \\
Inference decode, bs=32+ & 0.15--0.35 \\
Inference prefill & 0.30--0.50 \\
TinyML (TFLite Micro) & 0.05--0.15 \\
\bottomrule
\end{tabular}
\end{column}
\end{columns}

\vfill
\centering
\small\textcolor{gray}{You do not predict $\eta$ --- you measure it once and use it for what-if analysis.}
\end{frame}

% --- Slide 1.9: The Iron Law (Full) ---
\begin{frame}[shrink]{The Iron Law of ML Systems}
\note{[3 min] Walk through each denominator term. Point out that every
wall in the 22-wall taxonomy maps to exactly one term.
``Which term do you think is hardest to improve?''
% -- FLEX: [CORE]
}

\[
T = \frac{\text{FLOPs}}{N \;\times\; \text{Peak} \;\times\; \text{MFU}
  \;\times\; \eta_{\text{scaling}} \;\times\; \text{Goodput}}
\]

\vspace{0.3cm}
\small
\begin{tabular}{llll}
\toprule
\textbf{Term} & \textbf{Meaning} & \textbf{Reduced by} & \textbf{Walls} \\
\midrule
$N$ & Number of devices & Budget & --- \\
Peak & Raw hardware speed & GPU generation & 1 (Compute) \\
MFU & Software efficiency & FlashAttention, fusion & 2--3 \\
$\eta_{\text{scaling}}$ & Communication loss & BW, gradient compression & 14--16 \\
Goodput & Failure overhead & Checkpointing, FT & 15, 19 \\
\bottomrule
\end{tabular}

\vspace{0.3cm}
\centering
\textit{Every wall in the taxonomy attacks one of these five terms.}
\end{frame}

% --- Slide 1.10: Arithmetic Intensity Deep Dive ---
\begin{frame}[shrink]{Arithmetic Intensity: The Dial You Control}
\note{[2 min] ``Batch size is the primary knob. Each additional sample
in the batch reuses the same weights that are already loaded.
The compute grows linearly but memory stays constant.''
% -- FLEX: [CORE]
}

\small
\[
  \text{AI} = \frac{\text{FLOPs}}{\text{Bytes}} \approx
  \frac{2 \times \text{Params} \times B}{
    \underbrace{\text{Params} \times \text{bpp}}_{\text{weights}} +
    \underbrace{\text{Activations}(B)}_{\text{grows with } B}}
\]

\vspace{0.2cm}
\textbf{The batch-size knob:}
\begin{itemize}
  \item At $B=1$: AI $\approx$ 1 FLOP/byte $\Rightarrow$ \textbf{memory-bound}
  \item At $B=32$: AI $\approx$ 32 FLOP/byte $\Rightarrow$ approaching ridge
  \item At $B=256$: AI $\gg$ ridge $\Rightarrow$ \textbf{compute-bound}
\end{itemize}

\vspace{0.2cm}
\centering
\begin{tikzpicture}[scale=0.65]
  \draw[->, thick] (0,0) -- (7,0) node[right, font=\scriptsize] {Batch size};
  \draw[->, thick] (0,0) -- (0,4) node[above, font=\scriptsize] {Throughput};
  \draw[very thick, datastroke] (0.3,0.3) -- (3,3);
  \draw[very thick, computestroke] (3,3) -- (6.5,3);
  \fill[crimson] (3,3) circle (3pt);
  \node[above, font=\scriptsize, crimson] at (3,3.1) {Ridge};
  \node[below, font=\scriptsize, datastroke] at (1.3,0) {BW-bound};
  \node[below, font=\scriptsize, computestroke] at (5,0) {Compute-bound};
\end{tikzpicture}
\end{frame}

% --- Slide 1.11: Live Demo --- Batch Size Sweep ---
\begin{frame}[fragile,shrink=10]{Live Demo: Finding the Ridge Point}
\note{[3 min] Run this loop live. Show how the bottleneck flips
from Memory to Compute as batch size increases.
Before running, ask: ``At what batch size do you
predict the bottleneck will flip?'' Take guesses.
% -- FLEX: [CORE]
}

\small
\begin{lstlisting}
llama = mlsysim.Models.Language.Llama3_8B
hw    = mlsysim.Hardware.Cloud.H100

for bs in [1, 4, 16, 64, 128, 256]:
    p = mlsysim.Engine.solve(llama, hw, batch_size=bs)
    print(f"bs={bs:>3d}  {p.bottleneck:<8s}  "
          f"MFU={p.mfu:.3f}")
\end{lstlisting}

\vspace{0.3cm}
\begin{exampleblock}{What to observe}
  \begin{itemize}
    \item Bottleneck flips from \texttt{Memory} to \texttt{Compute}
    \item MFU climbs as batch size increases (better hardware utilization)
    \item Latency grows but throughput (tokens/s) improves
  \end{itemize}
\end{exampleblock}
\end{frame}

% --- Slide 1.12: Exercise 1 ---
\begin{frame}[fragile,shrink=10]{Exercise 1: Find the Crossover}
\note{[5 min] Attendees work individually.
At what batch size does Llama-3 8B on H100 transition
from memory-bound to compute-bound?
Expected: around bs=32--64 depending on precision.
Turn to your neighbor: did you get the same answer? Why or why not? 60 seconds.
% -- FLEX: [CORE] --- this is a critical hands-on moment.
}

\centering
\Large\bfseries
Hands-On Exercise\\[0.5cm]
\normalsize

\textbf{Question:} At what batch size does Llama-3 8B on H100
transition from memory-bound to compute-bound?

\vspace{0.3cm}
\begin{lstlisting}
for bs in range(1, 129):
    p = mlsysim.Engine.solve(llama, hw, batch_size=bs)
    if p.bottleneck == "Compute":
        print(f"Crossover at batch size {bs}")
        break
\end{lstlisting}

\vspace{0.3cm}
\textit{Bonus: Try the same on A100. Does the crossover happen
at the same batch size? Why or why not?}
\end{frame}

% --- Slide 1.13: The Ridge Point Explained ---
\begin{frame}[shrink]{The Ridge Point: Hardware DNA}
\note{[2 min] ``The ridge point is a property of the hardware,
not the workload. It tells you how many FLOPs per byte the chip
can sustain before compute becomes the ceiling.''
% -- FLEX: [CORE]
}

\small
\[
  \text{Ridge Point} = \frac{\text{Peak FLOPS}}{\text{Peak BW}}
  \;\;\bigl[\text{FLOP/byte}\bigr]
\]

\vspace{0.2cm}
\scriptsize
\begin{tabular}{llccc}
\toprule
\textbf{Vendor} & \textbf{Hardware} & \textbf{Peak FP16} & \textbf{HBM BW} & \textbf{Ridge} \\
\midrule
NVIDIA & H100 SXM  & 989 TFLOPS  & 3.35 TB/s & 295 FLOP/B \\
NVIDIA & B200      & 2.25 PFLOPS & 8.0 TB/s  & 281 FLOP/B \\
AMD    & MI300X    & 1,307 TFLOPS & 5.3 TB/s & 247 FLOP/B \\
Intel  & Gaudi\,3  & 1,835 TFLOPS & 3.7 TB/s & 496 FLOP/B \\
\bottomrule
\end{tabular}
\small

\vspace{0.3cm}
\begin{itemize}
  \item Higher ridge $\Rightarrow$ more workloads are memory-bound on this chip
  \item \alert{FLOPS grow faster than bandwidth across GPU generations}
  \item The memory wall is getting \textbf{worse}, not better
\end{itemize}
\end{frame}

% --- Slide 1.14: Predict Before You Peek #2 ---
\begin{frame}[shrink]{Predict: ResNet-50 vs Llama-3 8B}
\note{[2 min] ``ResNet-50 at batch 256 vs Llama-3 8B at batch 1.
Which is compute-bound and which is memory-bound?''
Give 30 seconds. Expected: ResNet at high batch is compute-bound;
Llama at bs=1 is memory-bound.
Turn to your neighbor: did you get the same answer? Why or why not? 60 seconds.
% -- FLEX: [CORE]
}

\centering
\Large\bfseries
Two workloads on the same H100.\\[0.3cm]
\normalsize
Which is compute-bound? Which is memory-bound?\\[0.5cm]

\pause

\small
\begin{tabular}{lcc}
\toprule
              & \textbf{ResNet-50 (bs=256)} & \textbf{Llama-3 8B (bs=1)} \\
\midrule
Total FLOPs   & $2.05 \times 10^{12}$       & $1.6 \times 10^{10}$ \\
Weight bytes   & 50 MB (FP16)               & 16 GB (FP16) \\
AI (FLOP/B)   & $\sim$41{,}000              & $\sim$1 \\
\midrule
\textbf{Regime} & \colorbox{computeblue}{Compute-bound} & \colorbox{datagreen}{Memory-bound} \\
\bottomrule
\end{tabular}

\vspace{0.3cm}
\pause
\alert{Same hardware, completely different bottlenecks.}\\
\textit{The workload determines the regime, not the GPU.}
\end{frame}

% --- Slide 1.15: Key Takeaway (Part 1) ---
\begin{frame}[shrink]{Part 1: Key Takeaway}
\note{[1 min] Summarize in one sentence. Repeat it twice.
``The bottleneck determines the speedup. Know your regime.''
% -- FLEX: [CORE]
}

\centering\Large

\textbf{The bottleneck determines the speedup.}\\[0.5cm]

\normalsize
\begin{itemize}
  \item The Roofline model tells you \textit{which} constraint is binding.
  \item Batch size is the primary knob that moves you between regimes.
  \item More FLOPS only helps if you are compute-bound.
  \item More bandwidth only helps if you are memory-bound.
  \item \texttt{Engine.solve()} answers this in one line.
\end{itemize}
\end{frame}

% --- Roadmap: After Break ---
\begin{frame}[shrink]{Roadmap: You Are Here}
\note{[1 min] Quick orientation after break. We now move from single-op analysis to serving.}

\centering\small
\begin{tabular}{rll}
\toprule
\textbf{Time} & \textbf{Part} & \textbf{Status} \\
\midrule
9:00--9:30   & Part 0: Welcome \& Setup     & \checkmark \\
9:30--10:30  & Part 1: Iron Law \& Roofline  & \checkmark \\
\rowcolor{crimson!12}
10:45--11:45 & \textbf{Part 2: Memory Walls \& Serving} & \textbf{$\leftarrow$ You are here} \\
11:45--12:00 & Part 3: Compression & \\
\bottomrule
\end{tabular}
\end{frame}

% =============================================================================
% PART 2: MEMORY WALLS & SERVING (12 slides)
% =============================================================================


\end{document}
